\section{Hitting-Set se encuentra en NP}

\subsection{Validador}

Para demostrar que el problema está en NP, proponemos un validador que, dada una instancia del problema ($A$, los $B_i$, y
$k$) y un certificado (el conjunto $C$ propuesto), verifique en tiempo polinomial si dicho certificado es, en efecto, una
solución válida.

A: conjunto de nn elementos sobre el cual se define el problema.

B: familia de mm subconjuntos B1,B2,…,BmB1​,B2​,…,Bm​, donde cada Bi⊆ABi​⊆A.

k: número entero que representa la cota máxima de tamaño permitida para el conjunto solución.

C: un subconjunto de A que se postula como solución al problema.

\begin{lstlisting}[language=Python]
def validadorHittingSet(A, B, k, C):
    if len(C) > k:
        return False

    for v in C:
        if v not in A:
            return False

    for Bi in B:
        if len(C.intersection(set(Bi))) == 0:
            return False

    return True
\end{lstlisting}

El validador hace lo siguiente:
\begin{itemize}
    \item Verifica que el tamaño del conjunto propuesto no supere a $k$. Esto es $\mathcal{O}(1)$.
    \item Verifica que todos los elementos propuestos pertenezcan efectivamente a $A$. Esto es $\mathcal{O}(k)$
    (asumiendo que if v not in A es $\mathcal{O}(1)$ al usar sets/diccionarios).
    \item Por cada uno de los $m$ subconjuntos $B_i$, verifica que $C$ tenga intersección no vacía con dicho subconjunto.
    Calcular esa intersección es, en el peor caso, $\mathcal{O}(|B_i|)$. Sumando para todos los $B_i$, esto es a lo sumo
    $\mathcal{O}(\sum_i |B_i|)$, que es, como mucho, $\mathcal{O}(m \cdot n)$ (cota burda, pero suficiente), y que en
    cualquier caso está acotado por el tamaño de la entrada (ya que los $B_i$ forman parte de la instancia).
\end{itemize}

Es decir, el validador corre en tiempo polinomial respecto del tamaño de la entrada (que ya incluye, de por sí, a los $B_i$
con sus respectivos elementos). Por lo tanto, Hitting-Set es un problema que se encuentra en NP.
